\scanpage{8}
\uline{Thm 1} Soit $\pi$ groupe discret, $E$ revêtement $\pi$-principal d'un espace $B$.
Considérons pour tout espace $X$ l'application
% Second pass: "tout" is written in blue ink, in the gap between "pour" and "espace" (Hua batch-01, p. 7:
% added in blue above the line).
\[
\operatorname{hot}(X,B)\xrightarrow{h_X}H^1(X,\pi)
\quad\bigl(\overset{\mathrm{df}}=\text{classes d'isom. de $\pi$-rev. princi-paux de $X$}\bigr)
\]
\figure{\begin{tikzpicture}
\draw[gray] (0,0) .. controls (.3,.35) and (.5,-.3) .. (.8,0)
.. controls (1.1,.35) and (1.3,-.2) .. (1.6,.05)
.. controls (1.9,.3) and (2.1,-.25) .. (2.4,0)
.. controls (2.7,.2) and (2.9,-.35) .. (3.2,-.1)
.. controls (3.5,.1) and (4.2,-.1) .. (4.6,-.08);
\end{tikzpicture}}
fonctorielle en $X$.
\margin{\unsure{5}}
% Second pass: the sign read 5 is a blue S-shaped stroke below H^1(X,\pi), under the wavy line, not in the
% margin; Hua batch-01 describes it as "un signe en forme de S".

Conditions équivalentes

1) $h$ est un isomorphisme de foncteurs,

\hspace*{2em}i.e. bijectif pour $\forall\ X$
% Second pass: "bijectif" is in blue ink; its first letters are written over another start of a word, with a
% stroke below (Hua batch-01: "bi" over a struck "in(jectif)"); the underlying letters are not legible here.

(On dit alors que $(E,B,\pi)$ est un $\pi$-rev. universel, $B$ s'appelle ``l'espace classifiant'' pour $\pi$)

2) $E$ est contractile

\mspage{7}
